% =============================================================================
%  SECCIONES REDACTADAS — PARTE 2
%  Capítulo 4 (Metodología): secciones 4.1, 4.5, 4.6, 4.7, 4.8
%  Incluye además correcciones a 4.2.1 y 4.4.3
% =============================================================================


% ╔═══════════════════════════════════════════════════════════════════════════╗
% ║  SUSTITUYE LA SECCIÓN 4.1                                                 ║
% ╚═══════════════════════════════════════════════════════════════════════════╝

\section{Arquitectura del sistema}

El sistema se organiza en tres capas desacopladas, según se recoge en la
Figura~\ref{fig:arquitectura}. Esta separación permite desarrollar y validar cada
una de forma independiente, y resulta especialmente conveniente en un trabajo con
componente de despliegue: la capa de aplicación puede evolucionar sin afectar al
procedimiento de entrenamiento, y viceversa.

\begin{figure}[H]
\centering
% NOTA: sustituir por el diagrama elaborado con draw.io o Excalidraw
\framebox[0.9\textwidth][c]{\rule{0pt}{5cm}\textit{[Diagrama de arquitectura]}}
\caption{Arquitectura general del sistema propuesto}
\label{fig:arquitectura}
\end{figure}

\subsection*{Capa de datos}

Responsable de la obtención, el tratamiento y el almacenamiento de la información.
Integra las tres fuentes descritas en la sección anterior mediante llamadas a sus
respectivas interfaces de programación, aplica el procedimiento de limpieza e
imputación y persiste el resultado en formato Parquet. Este almacenamiento
intermedio evita la dependencia de la red en ejecuciones sucesivas y garantiza la
reproducibilidad del análisis.

\subsection*{Capa de modelos}

Comprende la construcción de variables, el entrenamiento y la evaluación. Su
componente central es el módulo \texttt{prediccion.py}, que implementa la
generación de características y la función de predicción.

Conviene destacar una decisión de diseño con implicaciones prácticas
considerables: \textbf{este módulo es compartido por el procedimiento de
entrenamiento y por la aplicación}. La alternativa —implementar la construcción de
variables por duplicado— expone el sistema al fenómeno conocido como
\textit{training/serving skew}: si ambas implementaciones divergen, aunque sea en
un detalle menor, el modelo recibe en explotación entradas distintas de aquellas
sobre las que fue entrenado, y su rendimiento se degrada sin que ningún
procedimiento falle de manera visible. Compartir el código elimina esa
posibilidad por construcción.

\subsection*{Capa de aplicación}

Comprende la interfaz web, la obtención de datos en tiempo real y la generación de
recomendaciones. Se apoya en el modelo serializado que produce la capa anterior y
opera de forma autónoma respecto a ella.

Un principio rector de esta capa es la \textbf{degradación controlada}: el fallo
de una fuente externa no debe interrumpir el servicio. Si la previsión
meteorológica no está disponible, se recurre a una estimación climatológica y se
advierte de ello al usuario; si falla la consulta de precios, la aplicación
continúa operando sin esa información. Únicamente la serie de demanda resulta
imprescindible, y su ausencia se comunica mediante un mensaje explícito.


% ╔═══════════════════════════════════════════════════════════════════════════╗
% ║  SUSTITUYE LA SECCIÓN 4.5                                                 ║
% ╚═══════════════════════════════════════════════════════════════════════════╝

\section{Estrategia de validación}
\label{sec:validacion}

\subsection*{Improcedencia de la validación cruzada aleatoria}

La validación cruzada por particiones aleatorias, procedimiento estándar en
aprendizaje supervisado, resulta \textbf{metodológicamente inadmisible} en series
temporales. Al asignar observaciones a los conjuntos de entrenamiento y prueba sin
atender a su orden, el modelo se entrena con datos posteriores a aquellos que debe
predecir. Las métricas resultantes son optimistas y carecen de correspondencia con
el comportamiento en explotación.

\subsection*{Validación de origen rodante}

Se adopta en su lugar la validación de origen rodante
(\textit{rolling origin} o \textit{walk-forward}). El procedimiento fija un
instante de corte, entrena con la totalidad de las observaciones anteriores,
predice las veinticuatro horas siguientes y repite la operación desplazando el
corte:

\begin{verbatim}
Origen 1:  [----- entrenamiento -----][test]
Origen 2:  [------- entrenamiento -------][test]
Origen 3:  [--------- entrenamiento ---------][test]
\end{verbatim}

Este esquema reproduce las condiciones de explotación: en cada origen el modelo
dispone exclusivamente de información anterior al instante de predicción.

\subsection*{Elección del paso entre orígenes}

La separación entre orígenes exige una precaución que no siempre se explicita en
la literatura. Una primera configuración empleó un paso de ciento sesenta y ocho
horas —exactamente una semana— con el propósito de cubrir un periodo amplio. El
resultado fue que \textbf{los catorce orígenes recayeron sobre el mismo día de la
semana}, de modo que la evaluación excluía por completo sábados y domingos.

El sesgo introducido no era neutral. Beneficiaba específicamente a los modelos que
predicen a partir del día anterior, pues todos los días evaluados eran laborables
precedidos de laborables. Detectada la anomalía, se sustituyó el paso por ciento
veinte horas —cinco días—, valor que no es múltiplo de la semana y que con catorce
orígenes proporciona exactamente dos observaciones de cada día. Al repetir la
evaluación, el error del modelo ingenuo diario se elevó de forma apreciable, lo
que confirmó el diagnóstico.

La configuración definitiva es por tanto:

\begin{itemize}
  \item Horizonte de predicción: 24 horas
  \item Número de orígenes: 14
  \item Separación entre orígenes: 120 horas (5 días)
  \item Periodo cubierto: aproximadamente 70 días
\end{itemize}

\subsection*{Evaluación complementaria en días festivos}

El esquema anterior distribuye los orígenes de forma regular, por lo que la
coincidencia con un día festivo queda librada al azar. En la evaluación realizada
ningún origen recayó sobre una festividad, circunstancia que habría dejado sin
contrastar el hallazgo principal del análisis exploratorio.

Se construyó en consecuencia un conjunto de evaluación \textbf{dirigido},
seleccionando ocho orígenes de modo que el día completo a predecir fuese festivo.
Este conjunto está deliberadamente sesgado hacia casos difíciles y no representa
el rendimiento medio del sistema; se reporta como análisis complementario en la
sección~\ref{sec:festivos}.


% ╔═══════════════════════════════════════════════════════════════════════════╗
% ║  SUSTITUYE LA SECCIÓN 4.6 COMPLETA (4.6, 4.6.1, 4.6.2, 4.6.3)             ║
% ╚═══════════════════════════════════════════════════════════════════════════╝

\section{Modelos desarrollados}

Se han evaluado diez formulaciones, organizadas en tres niveles de complejidad
creciente. Todas comparten el mismo esquema de validación, los mismos orígenes y
las mismas métricas, de modo que la comparación resulta homogénea.

\subsection{Línea base}

\subsubsection*{Modelos ingenuos}

Antes de cualquier formulación estadística se establece el error que se obtiene
sin modelo alguno. Se consideran tres alternativas:

\begin{enumerate}
  \item \textbf{Ingenuo diario}: la demanda de cada hora del día siguiente
        coincide con la observada a esa misma hora del día actual.
  \item \textbf{Ingenuo semanal}: coincide con la del mismo día de la semana
        anterior. A diferencia del anterior, recoge el efecto laborable/festivo.
  \item \textbf{Perfil semanal}: media por hora del día sobre las cuatro últimas
        semanas.
\end{enumerate}

La inclusión de estas referencias responde a una observación recurrente en la
literatura \parencite{hong2016probabilistic}: un número considerable de trabajos
publicados presenta mejoras que en realidad no superan la predicción ingenua
estacional.

\subsubsection*{Modelos estadísticos clásicos}

\textbf{SARIMA.} Se ajusta un modelo SARIMA$(2,0,1)(1,1,1)_{24}$, con órdenes
orientados por las funciones de autocorrelación. El elevado coste computacional
del ajuste con periodo estacional 24 obliga a restringir el entrenamiento a una
ventana de sesenta días.

\textbf{SARIMAX con términos de Fourier.} Alternativa que representa las
estacionalidades diaria, semanal y anual mediante pares seno-coseno introducidos
como regresores externos, reservando la parte ARIMA para la autocorrelación de
corto plazo. La comparación de tiempos de ajuste justifica empíricamente esta
elección: mientras el SARIMA clásico requiere del orden de minutos sobre una
ventana reducida, la formulación con Fourier procesa la serie completa en
segundos.

Una primera especificación de este modelo empleó $d=0$ sin término constante. La
inspección de los resultados reveló un sesgo sistemático —subestimación en los
periodos de demanda elevada y sobrestimación en los de demanda reducida—
característico de la reversión a la media del conjunto de entrenamiento. La
corrección consistió en aplicar una diferenciación de primer orden, con lo que el
modelo opera sobre incrementos y se adapta al nivel local.

\textbf{Prophet.} Se configura con estacionalidad diaria, semanal y anual, e
incorpora el calendario de festivos con ventanas de un día a cada lado para
recoger el efecto de arrastre.

\subsection{Modelos de aprendizaje automático}

Se evalúan cuatro formulaciones sobre el conjunto de variables descrito en la
sección~\ref{sec:features}:

\begin{itemize}
  \item \textbf{Ridge}. Regresión lineal con regularización $L_2$, precedida de
        estandarización. Actúa como referencia lineal y permite aislar la
        aportación de la no linealidad.

  \item \textbf{Random Forest} \parencite{breiman2001rf}. Cien árboles, profundidad
        máxima doce y diez observaciones mínimas por hoja. La configuración es
        deliberadamente contenida: el procedimiento de validación exige catorce
        reentrenamientos sobre más de veinte mil observaciones, y una
        parametrización más ambiciosa resultaba impracticable.

  \item \textbf{XGBoost} \parencite{chen2016xgboost}. Cuatrocientos árboles, tasa
        de aprendizaje 0,05, profundidad máxima siete, submuestreo de
        observaciones y variables al 80\,\%.

  \item \textbf{LightGBM} \parencite{ke2017lightgbm}. Parámetros equivalentes,
        con cuarenta y ocho hojas por árbol.
\end{itemize}

\subsection{Sobre la optimización de hiperparámetros}

No se ha realizado una búsqueda sistemática de hiperparámetros. La decisión es
deliberada y se fundamenta en dos consideraciones.

En primer lugar, los resultados obtenidos con configuraciones convencionales
superan holgadamente los umbrales fijados en los objetivos del trabajo, por lo que
el margen de mejora esperable no justifica el coste computacional de una búsqueda
exhaustiva.

En segundo lugar, y con mayor peso, la evidencia acumulada en las competiciones M
\parencite{makridakis2022m5} indica que la calidad de las variables construidas
determina el rendimiento en mayor medida que el ajuste fino de los
hiperparámetros. El esfuerzo se ha concentrado en consecuencia sobre la ingeniería
de características, guiada por el análisis exploratorio.

La optimización mediante Optuna se recoge como línea de trabajo futuro.


% ╔═══════════════════════════════════════════════════════════════════════════╗
% ║  SUSTITUYE LA SECCIÓN 4.7 (conservando las ecuaciones existentes)         ║
% ╚═══════════════════════════════════════════════════════════════════════════╝

\section{Métricas de evaluación}

Se emplean cuatro métricas complementarias, ninguna de las cuales resulta
suficiente por sí sola.

El \textbf{error absoluto medio} expresa la desviación en megavatios y resulta
directamente interpretable, si bien depende de la escala:

\begin{equation}
\text{MAE} = \frac{1}{n}\sum_{i=1}^{n}\left|y_i - \hat{y}_i\right|
\end{equation}

La \textbf{raíz del error cuadrático medio} penaliza en mayor medida las
desviaciones grandes, lo que la hace sensible a los valores extremos:

\begin{equation}
\text{RMSE} = \sqrt{\frac{1}{n}\sum_{i=1}^{n}\left(y_i - \hat{y}_i\right)^2}
\end{equation}

El \textbf{error porcentual absoluto medio} es adimensional y constituye la
referencia habitual en la literatura sobre predicción de carga, lo que permite
contrastar los resultados con los trabajos publicados:

\begin{equation}
\text{MAPE} = \frac{100}{n}\sum_{i=1}^{n}\left|\frac{y_i - \hat{y}_i}{y_i}\right|
\end{equation}

El \textbf{error porcentual absoluto medio simétrico} corrige la asimetría del
anterior frente a sobrestimaciones e infraestimaciones:

\begin{equation}
\text{sMAPE} = \frac{100}{n}\sum_{i=1}^{n}
  \frac{2\left|y_i - \hat{y}_i\right|}{\left|y_i\right| + \left|\hat{y}_i\right|}
\end{equation}

\subsection*{Media y mediana}

Las métricas se calculan sobre cada uno de los catorce orígenes y se resumen
mediante media y mediana. La distinción no es accesoria: un único día atípico
—una festividad, un episodio meteorológico extremo— puede duplicar la media sin
que ello refleje el comportamiento habitual del modelo.

Se reportan por tanto ambos estadísticos, junto con la desviación típica entre
orígenes y el peor caso observado. La mediana constituye el criterio principal de
ordenación, por su robustez frente a valores extremos.

\subsection*{Métricas de entrenamiento y de producción}

Siguiendo la indicación expresa de la dirección del trabajo, se registran
conjuntamente el error dentro de muestra —sobre las observaciones de
entrenamiento— y fuera de muestra —sobre las predicciones efectivas—. Una
diferencia sustancial entre ambos indica que el modelo memoriza en lugar de
generalizar.

Debe precisarse un matiz relativo a los modelos estadísticos: en ellos el ajuste
dentro de muestra corresponde a una predicción de un paso adelante, mientras que
la evaluación se realiza a veinticuatro pasos. La comparación resulta en
consecuencia poco informativa. En los modelos de aprendizaje automático, en
cambio, ambos errores se calculan sobre el mismo horizonte y la diferencia sí es
interpretable.


% ╔═══════════════════════════════════════════════════════════════════════════╗
% ║  SUSTITUYE LA SECCIÓN 4.8                                                 ║
% ╚═══════════════════════════════════════════════════════════════════════════╝

\section{Desarrollo de la aplicación web}

\subsection*{Diseño funcional}

La aplicación se dirige a un usuario sin formación técnica en el ámbito
energético. Su función es traducir la predicción numérica en una recomendación
comprensible. Se estructura en cinco bloques:

\begin{enumerate}
  \item \textbf{Indicadores de situación}: demanda actual, máximo y mínimo
        previstos, y franja horaria más favorable.
  \item \textbf{Recomendación en lenguaje natural}, generada automáticamente a
        partir de la predicción.
  \item \textbf{Representación gráfica} de la serie histórica y la previsión, con
        las franjas recomendadas destacadas.
  \item \textbf{Tabla de horas} ordenada por demanda prevista, con el precio de
        mercado cuando está disponible.
  \item \textbf{Contexto meteorológico} y nota metodológica sobre el modelo y sus
        limitaciones.
\end{enumerate}

\subsection*{Tecnologías empleadas}

La interfaz se ha desarrollado con Streamlit \parencite{streamlit}, entorno que
permite construir aplicaciones de datos en Python sin requerir desarrollo
específico de interfaz. Las representaciones gráficas emplean Plotly, y el modelo
se serializa mediante \texttt{joblib}.

\subsection*{Obtención de datos en tiempo real}

La aplicación consulta tres fuentes en cada ejecución. La demanda reciente se
obtiene de la interfaz REData, en bloques de cinco días: las peticiones de mayor
amplitud provocan errores de servidor, dado que el servicio publica con resolución
de diez minutos. La previsión meteorológica procede del servicio de predicción
horaria municipal de AEMET, con horizonte de cuarenta y ocho horas. El precio de
mercado se recupera igualmente de REData, con carácter complementario.

\subsection*{Predicción más allá del horizonte nativo}

El modelo se ha entrenado para un horizonte de veinticuatro horas, límite
impuesto por la disponibilidad de los desfases: predecir el instante $t+25$
requeriría conocer $t+1$, aún no observado.

Para alcanzar las cuarenta y ocho horas se aplica un \textbf{esquema recursivo}:
las predicciones del primer tramo se incorporan a la serie y se emplean para
calcular los desfases del segundo. El procedimiento acumula error, circunstancia
que se cuantifica en el capítulo de resultados y de la que se advierte al usuario
en la propia interfaz.

\subsection*{Estrategia de despliegue}

La aplicación se despliega en Streamlit Community Cloud, servicio gratuito
integrado con el repositorio de código. El modelo serializado se versiona junto
con el resto del proyecto, mientras que las credenciales de AEMET se gestionan
mediante variables de entorno y quedan excluidas del control de versiones.


% ╔═══════════════════════════════════════════════════════════════════════════╗
% ║  CORRECCIÓN — último párrafo de la sección 4.2.1                          ║
% ║  El texto actual dice "mes a mes"; la implementación final trocea en       ║
% ║  bloques de cinco días para evitar errores del servidor                    ║
% ╚═══════════════════════════════════════════════════════════════════════════╝

% La descarga se ejecuta de forma fraccionada en bloques de pocos días, con
% reintentos automáticos y un intervalo de espera entre peticiones. El
% fraccionamiento no obedece únicamente a la cortesía con el servicio: las
% peticiones que abarcan periodos extensos provocan errores \texttt{502 Bad
% Gateway}, dado que el widget devuelve del orden de ciento cuarenta registros
% diarios. Los resultados se almacenan localmente en formato Parquet para evitar
% dependencias de red en ejecuciones posteriores.


% ╔═══════════════════════════════════════════════════════════════════════════╗
% ║  SUSTITUYE LA SECCIÓN 4.4.3                                               ║
% ║  Añadir \label{sec:features} al inicio de la sección 4.4                   ║
% ╚═══════════════════════════════════════════════════════════════════════════╝

\subsection{Prevención de fuga de información}
\label{sec:fuga}

La construcción de variables a partir de la propia serie objetivo expone el
procedimiento a la fuga de información, que constituye el riesgo metodológico
principal en la aplicación de aprendizaje automático a series temporales.

\subsubsection*{Formulación del problema}

Para un horizonte de veinticuatro horas, la predicción del instante $t+24$ debe
elaborarse desde el instante $t$. En ese momento, la disponibilidad efectiva de
cada variable es la que recoge la Tabla~\ref{tab:disponibilidad}.

\begin{table}[H]
\centering
\caption{Disponibilidad de las variables en el instante de predicción}
\label{tab:disponibilidad}
\begin{tabular}{|l|l|c|}
\hline
\rowcolor{grisclaro}
\textbf{Variable en $t+24$} & \textbf{Equivale a} & \textbf{¿Conocida en $t$?} \\
\hline
\texttt{lag\_1h}   & $y_{t+23}$ & No \\
\hline
\texttt{lag\_2h}   & $y_{t+22}$ & No \\
\hline
\texttt{lag\_3h}   & $y_{t+21}$ & No \\
\hline
\texttt{lag\_24h}  & $y_{t}$ & Sí \\
\hline
\texttt{lag\_168h} & $y_{t-144}$ & Sí \\
\hline
\texttt{media\_24h}  & media de $t \ldots t+23$ & No \\
\hline
\texttt{media\_168h} & ventana que alcanza $t+23$ & No \\
\hline
Calendario y meteorología & conocidos de antemano & Sí \\
\hline
\end{tabular}
\end{table}

Merece atención el caso de los estadísticos móviles. Su construcción en el
análisis exploratorio empleó un desplazamiento de una posición, adecuado para
predicción a un paso. Para un horizonte de veinticuatro horas, sin embargo, la
ventana de \texttt{media\_24h} en el instante $t+24$ abarca de $t$ a $t+23$,
tramo en su mayor parte desconocido. \textbf{La contaminación no se limita por
tanto a los desfases de orden reducido}, extremo que no resulta evidente a simple
vista.

\subsubsection*{Criterio adoptado}

Se establece la regla siguiente: para un horizonte de $h$ periodos únicamente son
admisibles los desfases de orden $k \geq h$ y los estadísticos móviles cuya
ventana se cierre antes del instante de predicción. En consecuencia:

\begin{itemize}
  \item Se descartan \texttt{lag\_1h}, \texttt{lag\_2h} y \texttt{lag\_3h}.
  \item Se descartan los estadísticos móviles del análisis exploratorio y se
        recalculan desplazando la serie el horizonte completo.
  \item Se conservan los desfases de 24, 25, 48 y 168 horas, así como las
        variables de calendario y meteorológicas.
\end{itemize}

\subsubsection*{Cuantificación del efecto}

El efecto de la fuga no se ha asumido, sino medido. Entrenando el mismo modelo con
ambos conjuntos de variables sobre idénticos orígenes de validación se obtiene que
\textbf{la inclusión de las variables contaminadas reduce el error aparente en un
59\,\%}.

Dicho de otro modo: un modelo que incorporase los desfases de orden reducido
aparentaría un rendimiento sustancialmente superior al que alcanzaría en
explotación, donde esas variables sencillamente no están disponibles. La magnitud
del efecto justifica el tratamiento explícito que recibe esta cuestión.
